% Reference entries, S (AN-11).  Law 1: os/iword/app/face_a.c F10 (the
% offered name, a typed character replaces it; save_named: an existing file
% asks "Replace A:\NAME? No (Yes)", y replaces, any other key asks for a
% name again; write_file: flatten first, "ERROR: File I/O error", "ERROR:
% Disk Full--Press Any Key to Continue"), wp5.c iw_wp5_write (the 327-byte
% prefix of WordPerfect 5.1 and the body), M_SIZE (IW_SIZE_MENU,
% IW_SIZE_ATTRS), paint_status (the status line: the name, Block on, Doc Pg
% Ln Pos, POS with Caps Lock; iw_fmt_inches).  Gated byte for byte against
% the real WordPerfect 5.1 by test-iword-face and test-flair-iword (MEMO ==
% MINT02's prefix + the body).
% Search (F2, Shift-F2) and Setup (Shift-F1) are P3 (bead initech-qvgd,
% WL-0146): face_a.c function_key (IW_FK_search_forward / _backward,
% IW_FK_setup from spec/iword/fkeys.json), search_key (the string kept for
% the next Search and replaced by the first key typed; F2 / Shift-F2 / Esc
% search, F1 cancels with the cursor kept, Up / Down the direction, Left /
% Right / Backspace / Delete by whole codes, F6 / F8 an attribute's on or
% off code by the state at that point of the string, Tab [Tab], F4
% [->Indent], Enter [HRt], Ctrl-Enter [HPg]; 256 bytes, QUERY_CAP),
% os/iword/core/search.c iw_search (a lower-case letter finds either case,
% a capital only a capital; the cursor after the match forward, at its
% start backward; "* Not found *"), paint_setup / setup_key / settings_save
% (the Setup screens, the grey rows, Environment 3, Backup 1 / 2, Units 8
% with i c p w u; F7 one level up, F1 / Esc to the document; IWORD.SET,
% spec/iword/settings.h, eight bytes; the units and the backup choice are
% recorded and shown, nothing else reads them).
% GATES: test-pubs-iword test-iword-face test-flair-iword test-iword-wp5-roundtrip test-iword-codes test-iword-reveal-codes test-iword-search test-iword-setup test-flair-iword-search test-flair-iword-setup
\Gates{test-pubs-iword test-iword-face test-flair-iword test-iword-wp5-roundtrip test-iword-codes test-iword-reveal-codes test-iword-search test-iword-setup test-flair-iword-search test-flair-iword-setup}
\begin{Feature}{Save}
Save writes the document on the screen to the data volume as a file, and
leaves it on the screen so that you can go on working.  Save often: a
document that is only on the screen is lost if the computer is switched off.

To save a document,
\begin{WPSteps}
\item Press \FK{Save}{F10}.  The status line shows \Scr{Document to be
Saved:} and, if the document has a name, that name.
\PullDown{Select \Mi{S}ave from the \Mi{F}ile menu.}
\item To save under the name shown, press \K{Enter}.
\Or
Type a new name (the first character you type replaces the name shown) and
press \K{Enter}.
\item If a document with that name is already on the data volume, the status
line asks \Scr{Replace A:\wbs\textit{name}? No (Yes)}.  Type \Typed{y} to
replace it with the document on the screen.  Any other key asks you for a
name again, so that you can choose another.
\end{WPSteps}
\Documents{wpkey}{F10}
When the document is saved, its name appears at the left of the status line.

\BeginReplay{ref-save}
For example, type \TypeRun{Budget memo}, press \Press{Save}{F10}, type
\TypeName{budget} and press \PressKey{Enter}.  Press \Press{Save}{F10}
again: InitechWord offers the name.
\WordScreen[rows=1]{ref-save-1}{Save offers the document's name}
Press \PressKey{Enter}.  The file \Scr{BUDGET} exists, so InitechWord asks:
\WordScreen[rows=1]{ref-save-2}{Saving over a document that is already on the data volume}
Type \TypeName{y}: the document is saved again.
\EndReplay

\Notes
\Topic{File Format}
InitechWord saves documents in the format of WordPerfect 5.1, writing the
text and codes as WordPerfect 5.1 does (see \InRef{WordPerfect 5.1
Documents} and Appendix D).

\Topic{Messages}
\Scr{ERROR: Invalid file name} means the name breaks the rules for names
(see \InRef{Filenames}).  \Scr{ERROR: File I/O error} means the file could not
be made on the data volume.  \Scr{ERROR: Disk Full--Press Any Key to
Continue} means the data volume has no room for the document; the document
is still on the screen.  After each of these, press any key to return to
the name prompt.  \Scr{ERROR: Document too large} means the document is
larger than a file InitechWord can write (see Appendix H); press any key to
return to the document.

\Topic{Saving When You Exit}
Exit (F7) asks whether to save the document before you leave it, and saves
it as Save does (see \InRef{Exit}).

\SeeAlso{Exit; Filenames; Retrieve}
\end{Feature}

\begin{Feature}{Screen}
The editing screen is the InitechWord window as you type.  It shows 23
lines of the document, 76 characters across, and the status line under
them.  The first column of the screen is the left margin of a new document,
one inch from the left edge of the paper.

\Notes
\Topic{Colors}
Text is white on a blue field.  The cursor is the cell drawn in reverse,
blue on white.  A block is black on grey; so is the ruler of Reveal Codes.

\Topic{Faces}
Bold, italic and underlined text is drawn in that face.  Text with any other
appearance or size is drawn as ordinary text (see \InRef{Appearance} and
\InRef{Size}).

\Topic{Hard Pages and Spacing}
A hard page is drawn as a row of equal signs below the line it ends.  With
line spacing of more than one, blank rows are drawn under each line (see
\InRef{Line Spacing}).

\Topic{Margins Under One Inch}
The first column of the screen is always one inch from the edge of the
paper.  With a left margin of less than one inch, the start of each line
lies to the left of the window and is not shown; the status line's Pos still
tells where the cursor is.

\Topic{Scrolling}
When the cursor moves above the top line or below the bottom line, the
screen moves to keep it in view.

\SeeAlso{Reveal Codes; Status Line; Window}
\end{Feature}

\begin{Feature}{Search}
Search finds text, codes or both in the document.  It searches forward,
from the cursor toward the end of the document, or backward, from the
cursor toward the beginning, and stops at the first place where it finds
what you typed.

To search,
\begin{WPSteps}
\item Press \FK{\rarr Search}{F2} to search forward
\Or
Press \FK{\larr Search}{Shift-F2} to search backward.
\item The status line shows \Scr{-> Srch:} (forward) or \Scr{<- Srch:}
(backward), followed by what you last searched for, if anything.  Type the
text you want to find.  The first character you type replaces the text
shown.
\item Press \K{F2} to start the search.
\end{WPSteps}
\FlairShot{search}{640}{440}{The search string typed at the \Scr{-> Srch:} prompt}
\Documents{wpkey}{F2}
\Documents{wpkey}{Shift-F2}
\Documents{wpkey}{Srch:F2}
\Documents{wpkey}{Srch:Esc}
If InitechWord finds the text, the cursor stops just after it when you
search forward, and on its first character when you search backward.  If it
does not, the status line shows \Scr{* Not found *} and the cursor stays
where it was; press any key to go on.

\BeginReplay{ref-search}
For example, type
\Type{Every report needs a cover sheet.  The cover sheets are in the supply room.}
press \PressKey{Enter} and type
\Type{Ask Milton for Cover Sheets if the room is empty.}
Press \PressKey{Home, Home, Up} to go to the top of the document.  Press
\Press{\rarr Search}{F2} and type \TypeName{cover sheet}:
\WordScreen[rows=3]{ref-search-1}{The search string on the status line}
Press \PressKey{F2}.  The cursor stops after the first \textit{cover sheet}:
\WordScreen[rows=3]{ref-search-2}{The cursor after the text found}
Press \PressKey{F2} twice.  The first F2 shows \Scr{-> Srch:cover sheet}
again; the second searches for it, and the cursor stops after the
\textit{cover sheet} of \textit{cover sheets}.  Press \PressKey{F2} twice
more:\Keys{F2 F2}
\WordScreen[rows=3]{ref-search-3}{A string in lower case finds capitals as well}
Search found \textit{Cover Sheet}: a lower-case letter in the search string
finds the letter in either case.
\EndReplay

\Notes
\Topic{Block}
While Block is on, F2 and Shift-F2 do nothing.  Turn Block off to search.

\Topic{Canceling a Search}
Press \FK{Cancel}{F1} at the \Scr{Srch:} prompt to leave without searching.
The cursor stays where it was, and the search string is kept for the next
search.
\Documents{wpkey}{Srch:F1}

\Topic{Capital Letters}
A lower-case letter in the search string finds the same letter in upper or
lower case: \Typed{memo} finds \Scr{memo}, \Scr{Memo} and \Scr{MEMO}.  A
capital letter finds only a capital: \Typed{Memo} finds \Scr{Memo} and
\Scr{MEMO}, but not \Scr{memo}.

\Topic{Codes in the Search String}
To search for a code, press its key at the \Scr{Srch:} prompt.  The code is
put in the string as Reveal Codes shows it:
\par\addvspace{4pt}
\noindent\begin{tabular}{@{}ll@{}}
\textbf{Key} & \textbf{Code in the string}\\
\K{F6} & \Code{[BOLD]}, or \Code{[bold]} after a \Code{[BOLD]} not yet ended\\
\K{F8} & \Code{[UND]}, or \Code{[und]} after an \Code{[UND]} not yet ended\\
\K{Tab} & \Code{[Tab]}\\
\K{F4} & \Code{[\rarr Indent]}\\
\K{Enter} & \Code{[HRt]}\\
\K{Ctrl-Enter} & \Code{[HPg]}\\
\end{tabular}
\par\addvspace{4pt}
\Documents{wpkey}{Srch:F6}
\Documents{wpkey}{Srch:F8}
\Documents{wpkey}{Srch:Tab}
\Documents{wpkey}{Srch:F4}
\Documents{wpkey}{Srch:Enter}
\Documents{wpkey}{Srch:Ctrl-Enter}
A string can hold codes and text together: \Code{[BOLD]}\Typed{Note} finds
the word \textit{Note} only where bold begins just before it.  Codes of other
kinds cannot be put in the search string.

\BeginReplay{ref-search-code}
For example, type \TypeRun{Read the }, press \Press{Bold}{F6}, type
\TypeRun{first}, press \Press{Bold}{F6} again and type
\TypeRun{ page.  The first page is the summary.}  Press
\PressKey{Home, Home, Up}, press \Press{\rarr Search}{F2}, press
\PressKey{F6} and type \TypeName{first}:
\WordScreen[rows=1]{ref-search-code-1}{A code in the search string}
Press \PressKey{F2}.  The cursor stops after the bold \textit{first}, not
after the second one, which has no code before it.
\WordScreen[rows=1]{ref-search-code-2}{The bold word found}
\EndReplay

\Topic{Direction}
At the \Scr{Srch:} prompt, \K{Up} changes the search to backward
(\Scr{<- Srch:}) and \K{Down} to forward (\Scr{-> Srch:}).  \K{F2},
\K{Shift-F2} and \K{Esc} all start the search in the direction the prompt
shows.
\Documents{wpkey}{Srch:Up}
\Documents{wpkey}{Srch:Down}

\Topic{Editing the Search String}
\K{Left} and \K{Right} move the cursor through the search string, and
\K{Backspace} and \K{Delete} erase, as on the editing screen; each moves
past or erases a whole code at a time.  Once you have pressed one of them,
typing no longer replaces the string: what you type goes in at the cursor.
\Documents{wpkey}{Srch:Left}
\Documents{wpkey}{Srch:Right}
\Documents{wpkey}{Srch:Backspace}
\Documents{wpkey}{Srch:Delete}

\Topic{Keys That Act as Other Keys}
At the \Scr{Srch:} prompt InitechWord does not look at Shift, Alt or Ctrl on
these keys: each combination does what the key alone does.
\par\addvspace{4pt}
\noindent{\small\begin{tabular}{@{}lll@{}}
\WPPromptAlias{Srch:}{Shift-F1}{F1}, & \WPPromptAlias{Srch:}{Alt-F1}{F1}, & \WPPromptAlias{Srch:}{Ctrl-F1}{F1}\\
\WPPromptAlias{Srch:}{Shift-F2}{F2}, & \WPPromptAlias{Srch:}{Alt-F2}{F2}, & \WPPromptAlias{Srch:}{Ctrl-F2}{F2}\\
\WPPromptAlias{Srch:}{Shift-Esc}{Esc}, & \WPPromptAlias{Srch:}{Alt-Esc}{Esc}, & \WPPromptAlias{Srch:}{Ctrl-Esc}{Esc}\\
\WPPromptAlias{Srch:}{Shift-Up}{Up}, & \WPPromptAlias{Srch:}{Alt-Up}{Up}, & \WPPromptAlias{Srch:}{Ctrl-Up}{Up}\\
\WPPromptAlias{Srch:}{Shift-Down}{Down}, & \WPPromptAlias{Srch:}{Alt-Down}{Down}, & \WPPromptAlias{Srch:}{Ctrl-Down}{Down}\\
\WPPromptAlias{Srch:}{Shift-Left}{Left}, & \WPPromptAlias{Srch:}{Alt-Left}{Left}, & \WPPromptAlias{Srch:}{Ctrl-Left}{Left}\\
\WPPromptAlias{Srch:}{Shift-Right}{Right}, & \WPPromptAlias{Srch:}{Alt-Right}{Right}, & \WPPromptAlias{Srch:}{Ctrl-Right}{Right}\\
\WPPromptAlias{Srch:}{Shift-Backspace}{Backspace}, & \WPPromptAlias{Srch:}{Alt-Backspace}{Backspace}, & \WPPromptAlias{Srch:}{Ctrl-Backspace}{Backspace}\\
\WPPromptAlias{Srch:}{Shift-Delete}{Delete}, & \WPPromptAlias{Srch:}{Alt-Delete}{Delete}, & \WPPromptAlias{Srch:}{Ctrl-Delete}{Delete}\\
\WPPromptAlias{Srch:}{Shift-Tab}{Tab}, & \WPPromptAlias{Srch:}{Alt-Tab}{Tab}, & \WPPromptAlias{Srch:}{Ctrl-Tab}{Tab}\\
\WPPromptAlias{Srch:}{Shift-Enter}{Enter}, & \WPPromptAlias{Srch:}{Alt-Enter}{Enter} & \\
\end{tabular}}
\par\addvspace{4pt}
(Ctrl-Enter is the hard page code, as above.)

\Topic{Length of the String}
The search string holds up to 256 bytes of the document file: one for each
character, \Code{[HRt]} or \Code{[HPg]}, and the length given in Appendix D
for each other code.  Keys typed after the string is full are ignored.

\Topic{Repeating a Search}
InitechWord keeps the search string until you leave InitechWord.  To find the next
place, press F2 twice: once to show the string, once to search.

\Topic{The Document Is Not Changed}
Searching moves only the cursor.  It does not change the document, and a
document that was saved is still saved after a search.

\Topic{Words Within Words}
Search finds the string wherever it occurs, also inside a longer word:
\Typed{form} finds \Scr{form}, \Scr{forms} and \Scr{platform}.  To find a
whole word, put a space before or after it in the string.

\SeeAlso{Block; Codes; Cursor Movement; Replace}
\end{Feature}

\begin{Feature}{Setup}
Setup records two preferences of InitechWord: the units in which
measurements are to be shown, and whether a copy of a document is to be kept
when it is saved over.  InitechWord keeps them in a file on the data volume,
so that they are the same each time you start it.

To change a setting,
\begin{WPSteps}
\item Press \FK{Setup}{Shift-F1}.  The Setup screen appears.
\item Select \Ch{Environment}{3}.
\item Select \Ch{Backup Options}{1} or \Ch{Units of Measure}{8}, then the
setting, and give its value (see the Notes below).
\item Press \FK{Exit}{F7} to go back one screen at a time, or \FK{Cancel}{F1}
to return to the document at once.
\end{WPSteps}
\Documents{wpkey}{Shift-F1}
\Documents{wpmenu}{Shift-F1,3=Environment}
\Documents{wpmenu}{Shift-F1,3,1=Backup Options}
\Documents{wpmenu}{Shift-F1,3,8=Units of Measure}

\FlairShot{setup}{640}{440}{The Setup screen}
The choices printed in grey are WordPerfect settings InitechWord does not
have; typing their numbers does nothing.

\BeginReplay{ref-setup}
For example, press \Press{Setup}{Shift-F1} and select
\Choose{Environment}{3} and \Choose{Units of Measure}{8}:
\WordScreen[rows=13]{ref-setup-2}{Setup: Units of Measure}
Press \FK{Exit}{F7}\Keys{F7} twice\Keys{F7} and \FK{Cancel}{F1}\Keys{F1} to
return to the document.
\EndReplay

\Notes
\Topic{Backup Options}
Select \Ch{Backup Options}{1} on the Environment screen.  Timed Document
Backup (1) is printed in grey: InitechWord has none.  Select
\Ch{Original Document Backup}{2}; the status line asks
\Scr{Original Document Backup? No (Yes)}.  Type \Typed{y} or \Typed{n}.
The screen shows the new value.
\Documents{wpmenu}{Shift-F1,3,1,2=Original Document Backup}
InitechWord records the answer, but in this release it keeps no backup copy
either way: Save replaces a document as described under \InRef{Save}.

\Topic{Exit and Cancel}
On every Setup screen \FK{Exit}{F7} goes back one screen: from Units of
Measure or Backup Options to Environment, from Environment to Setup, and
from Setup to the document.  \FK{Cancel}{F1} and \K{Esc} return to the
document from any Setup screen.

\Topic{The Setup File}
InitechWord keeps the settings in the file \Scr{IWORD.SET} in the top level
of the data volume, and writes it each time you change a setting.  The file
is eight bytes long: the letters \Scr{IWS1}, the letter of the unit for
margins and tabs, the letter of the unit for the status line, 1 if Original
Document Backup is Yes (otherwise 0), and a 0.  If the file is missing or
not in this form, InitechWord starts with Inches and No.  If the file cannot
be written, the status line shows \Scr{ERROR: File I/O error} and the setting
is not changed.

\Topic{Units of Measure}
Select \Ch{Units of Measure}{8} on the Environment screen, then
\Ch{Display and Entry of Numbers for Margins, Tabs, etc.}{1} or
\Ch{Status Line Display}{2}.  The status line shows
\Scr{Units of Measure:}; type the letter of a unit:
\par\addvspace{4pt}
\noindent\begin{tabular}{@{}ll@{}}
\textbf{Letter} & \textbf{Unit}\\
\Typed{i} & Inches\\
\Typed{c} & Centimeters\\
\Typed{p} & Points\\
\Typed{w} & 1200ths of an inch\\
\Typed{u} & WordPerfect 4.2 units\\
\end{tabular}
\par\addvspace{4pt}
\Documents{wpmenu}{Shift-F1,3,8,1=Display and Entry of Numbers for Margins, Tabs, etc.}
\Documents{wpmenu}{Shift-F1,3,8,1,i=Inches}
\Documents{wpmenu}{Shift-F1,3,8,1,c=Centimeters}
\Documents{wpmenu}{Shift-F1,3,8,1,p=Points}
\Documents{wpmenu}{Shift-F1,3,8,1,w=1200ths of an inch}
\Documents{wpmenu}{Shift-F1,3,8,1,u=WordPerfect 4.2 units}
\Documents{wpmenu}{Shift-F1,3,8,2=Status Line Display}
\Documents{wpmenu}{Shift-F1,3,8,2,i=Inches}
\Documents{wpmenu}{Shift-F1,3,8,2,c=Centimeters}
\Documents{wpmenu}{Shift-F1,3,8,2,p=Points}
\Documents{wpmenu}{Shift-F1,3,8,2,w=1200ths of an inch}
\Documents{wpmenu}{Shift-F1,3,8,2,u=WordPerfect 4.2 units}
Any other key is ignored.  The Units of Measure screen shows the unit
chosen.  In this release InitechWord records the units but goes on showing
and taking every measurement in inches (see \InRef{Units of Measure}).

\SeeAlso{Save; Units of Measure}
\end{Feature}

\begin{Feature}{Size}
Size changes the size of text: superscript and subscript, and the sizes from
fine to extra large.  Like the appearances, each size is a pair of codes
around the text.

To give text a size as you type it,
\begin{WPSteps}
\item Press \FK{Font}{Ctrl-F8}, then select \Ch{Size}{1}.  The Size menu
appears on the status line:
\StatusLine{1 Suprscpt; 2 Subscpt; 3 Fine; 4 Small; 5 Large; 6 Vry Large; 7 Ext Large: 0}
\item Select the size by its number.
\item Type the text.
\item Press \KC{Ctrl-F8}{1} and select the same size again to move past its
off code
\Or
Press \FK{Font}{Ctrl-F8} and select \Ch{Normal}{3}.
\end{WPSteps}
\Documents{wpmenu}{Ctrl-F8,1=Size}
\Documents{wpmenu}{Ctrl-F8,1,1=Suprscpt}
\Documents{wpmenu}{Ctrl-F8,1,2=Subscpt}
\Documents{wpmenu}{Ctrl-F8,1,3=Fine}
\Documents{wpmenu}{Ctrl-F8,1,4=Small}
\Documents{wpmenu}{Ctrl-F8,1,5=Large}
\Documents{wpmenu}{Ctrl-F8,1,6=Vry Large}
\Documents{wpmenu}{Ctrl-F8,1,7=Ext Large}
\par\addvspace{6pt}
\noindent\begin{tabular}{@{}lll@{}}
\textbf{Selection} & \textbf{Size} & \textbf{Codes}\\
Suprscpt (1) & superscript & \Code{[SUPRSCPT]} \Code{[suprscpt]}\\
Subscpt (2) & subscript & \Code{[SUBSCPT]} \Code{[subscpt]}\\
Fine (3) & fine & \Code{[FINE]} \Code{[fine]}\\
Small (4) & small & \Code{[SMALL]} \Code{[small]}\\
Large (5) & large & \Code{[LARGE]} \Code{[large]}\\
Vry Large (6) & very large & \Code{[VRY LARGE]} \Code{[vry large]}\\
Ext Large (7) & extra large & \Code{[EXT LARGE]} \Code{[ext large]}\\
\end{tabular}
\par\addvspace{4pt}

\Notes
\Topic{On the Screen}
The editing screen draws text of every size in the same cells as ordinary
text.  Press Reveal Codes to see which size text has.

\Topic{Size and Appearance Together}
Text can have a size and appearances at once; each is its own pair of codes
(see \InRef{Appearance}).

\SeeAlso{Appearance; Block, Attributes; Font}
\end{Feature}

\begin{Feature}{Status Line}
The status line is the bottom line of the InitechWord window.  While you
edit, it shows the name of the document at its left and where the cursor is
at its right.  InitechWord also uses it for menus, prompts, questions and
messages.

To read the status line,
\begin{WPSteps}
\item Look at the left end: the document's name once it has one, as
\Scr{A:\wbs MEMO}; or \Scr{Block on} while Block is on; or
\Scr{Move cursor; press Enter to retrieve.} while moved text waits.
\item Look at the right end: \Scr{Doc 1 Pg 1 Ln 1" Pos 1"}.
\end{WPSteps}

\BeginReplay{ref-status}
For example, type
\Type{Each line of a single-spaced paragraph is one sixth of an inch below the line before it, so Ln grows by 0.17 inch a line.}
The cursor is on the second line, after the period:
\WordScreen[rows=2]{ref-status}{Ln and Pos at the end of a two-line paragraph}
\EndReplay

\Notes
\Topic{Doc}
The document number is always 1.

\Topic{Ln}
Ln is the distance from the top edge of the page to the line the cursor is
on, in inches: \Scr{1"} for the first line under the one inch top margin,
then a sixth of an inch for each single-spaced line (\Scr{1.17"},
\Scr{1.33"}, \ldots).  With other line spacings, each line adds its own
spacing (see \InRef{Line Spacing}).

\Topic{Measurements}
Ln and Pos are shown in inches, to two decimal places at most, without
trailing zeros: \Scr{1"}, \Scr{1.5"}, \Scr{2.17"}.

\Topic{Pg}
Pg is the page the cursor is on (see \InRef{Page Break, Hard}).

\Topic{Pos}
Pos is the distance from the left edge of the paper to the cursor, in
inches: \Scr{1"} at the standard left margin, then a tenth of an inch for
each character.  While Caps Lock is on, Pos is shown as \Scr{POS} (see
\InRef{Caps Lock}).

\SeeAlso{Caps Lock; Cursor Movement; Line Spacing; Page Break, Hard}
\end{Feature}
